\section{Conditions holonomiques du sélecteur gravitationnel général}

Le second schéma gravitationnel part de
\begin{equation}
\label{rm:eq:gravity-delta-theta}
\delta_\theta
=\frac{A_{\rm deg}-C^*_{\rm deg}}{360}
\end{equation}
et d'une longueur de route
\(\Lambda_{27}=\mathcal L_{\rm hol}(\widetilde{\mathcal K}_{27})\)
dans un relèvement enrichi du retour. L'expression
\begin{equation}
\label{rm:eq:gravity-holonomic-template}
\boxed{
G_{\rm hol}
=\frac{c_{\rm int}^3}{\hbar_{\rm int}}
\bigl(\delta_\theta\Lambda_{27}\bigr)^2}
\end{equation}
est dimensionnellement exacte. Sa promotion en sélecteur général requiert que
\(\mathcal L_{\rm hol}\) soit naturel sous changement cyclique de base,
conjugaison, réétiquetage et subdivision, et qu'il possède une normalisation
unique stable sous raffinement.

Il existe en outre un contrôle de torsion : tout caractère multiplicatif positif
\(\chi:K_{27}\to\R_+^\times\) sur un cycle observable fini d'ordre
quatre est trivial, car \(\chi(g)^4=1\) et \(\chi(g)>0\) impliquent
\(\chi(g)=1\). Une réponse holonomique non triviale doit donc
être définie sur le relèvement sans torsion ou sur la mémoire
enrichie ; elle ne peut pas
s'obtenir en renommant un caractère positif du cycle fini.

La carte
\begin{equation}
\label{rm:eq:gravity-decimal-chart}
10^{-11}\left(6+\frac{674}{1000}+\frac{30}{10^5}\right)
=6.67430\times10^{-11}
\end{equation}
est une évaluation exacte d'entrées déclarées. Elle ne fixe pas
\(\Lambda_{27}\) et ne démontre pas que \cref{rm:eq:gravity-holonomic-template}
prenne cette valeur.

\begin{keyresult}
La condition de périodicité détermine \(G\) sans employer son développement
décimal : la structure chronologique CT108 sélectionne \((54/\pi)^2\), la
droite gravitationnelle fixe le type dimensionnel et la famille bidirectionnelle
conserve \(G\hbar\). Le couplage obtenu est représenté sur le
support tensoriel construit par \(P_5\), tandis que la réduction
physique transverse sans trace s'effectue au moyen d'une application distincte de
rang deux.
\end{keyresult}

\begin{falsifier}
Le bloc circulaire est réfuté si
\(r_g(\theta)\neq\theta(\pi/54)\ell_0\), s'il existe une autre solution positive
de \cref{rm:eq:gravity-closing-condition}, ou si
\(G_a\hbar_a/c^3\neq\theta_{108}^{\rm circ}\ell_0^2\). Le bloc tensoriel est
réfuté si \(\pi_{\rm ico}\) ne préserve pas l'incidence ou n'entrelace pas l'
action déclarée, si son stabilisateur n'est pas \(C_5\), si \(P_5\) n'est pas
un idempotent autoadjoint de trace cinq, si
\(\Gamma_0\Gamma_0^*\neq(8/5)I\), ou si \(\Lambda(n)\) n'a pas pour rang deux.
La lecture de cinq polarisations est réfutée si les tenseurs de
\cref{rm:eq:gravity-helicity2-real,rm:eq:gravity-helicity1-real,rm:eq:gravity-helicity0-real}
ne forment pas une base orthonormée ou si \(\Lambda(n)\) retient un poids différent de
\(\pm2\).
Le sélecteur holonomique général est résolu par obstruction : sans une
\(\mathcal L_{\rm hol}\) naturelle, normalisée et stable, il n'existe pas de
sélection univoque compatible avec tous les réétiquetages ; utiliser
\cref{rm:eq:gravity-decimal-chart} pour la choisir entraîne une circularité et
invalide la construction.
\end{falsifier}

\begin{statusbox}
Transducteur, unicité du sélecteur sous la condition de périodicité CT108 et
famille de Planck : \status{Théorème interne exact sous la condition géométrique déclarée}.
La descente \(\Omega_{12}^{\rm ico}\simeq A_5/C_5\) est exacte sous les quatre
hypothèses explicites de \(\pi_{\rm ico}\) ; la construction HMT de cette
application appartient au corridor exceptionnel de provenance. La descente étant fixée,
projecteur \(P_5\), entrelaceur \(\Gamma_5\), base de cinq hélicités,
réduction TT et spectre de \(\mathbb G_5\) : \status{Théorème interne exact}.
Lecture massive, équations de champ et réalisation
Palatini--Cartan--Holst : \status{Réalisation physique typée}. Schéma
holonomique : \status{Démontrée en homogénéité}; l'inexistence d'un
sélecteur général univoque sans normalisation supplémentaire est une
\status{obstruction de canonicité démontrée}. Carta SI:
\status{Contraste métrologique externe}.
\end{statusbox}
